% =============================================================================
%  Thesis information
%  ---------------------------------------------------------------------------
%  Fill in the fields below. They feed the title page, the back cover, the
%  résumé/abstract pages and the PDF metadata. Set a field to {} to hide the
%  line it produces.
% =============================================================================

% ── Language ──────────────────────────────────────────────────────────────────
% French is the main language. For a thesis written in English, replace the
% line below with \thesisfrenchfalse: headings, lists, the title page and the
% bibliography all switch to English. The résumé and the abstract are always
% both included.
\thesisfrenchtrue

% ── Institution ───────────────────────────────────────────────────────────────
\newcommand{\University}{Université [Nom]}
\newcommand{\Faculty}{UFR / Faculté / École [Nom]}
\newcommand{\Department}{Département [Nom]}

% ── Degree ────────────────────────────────────────────────────────────────────
% e.g. Mémoire de Master 2 · Mémoire de recherche · Rapport de stage
\newcommand{\ThesisType}{Mémoire de Master}
% e.g. Master Informatique · Master Économie appliquée
\newcommand{\Degree}{Master [Mention]}
% e.g. Parcours Intelligence artificielle — {} if there is no track
\newcommand{\Track}{Parcours [Nom du parcours]}

% ── Title ─────────────────────────────────────────────────────────────────────
% The title page shows the title in the main language; the back cover shows
% both.
\newcommand{\TitleFR}{Titre du mémoire}
\newcommand{\TitleEN}{Thesis Title}
\newcommand{\Subtitle}{}

% ── Author ────────────────────────────────────────────────────────────────────
\newcommand{\Author}{Prénom NOM}

% ── Host organisation ─────────────────────────────────────────────────────────
% The research laboratory or the company where the work was carried out.
% Set both to {} if there is none.
\newcommand{\HostName}{[Laboratoire / Entreprise]}
\newcommand{\HostAddress}{[Adresse]}

% ── Logos ─────────────────────────────────────────────────────────────────────
% Upload the image file and give its path, e.g. {figures/logo-universite.png}.
% {} shows a dashed placeholder; {none} removes the logo slot.
\newcommand{\UniversityLogo}{}
\newcommand{\HostLogo}{}

% ── Defence ───────────────────────────────────────────────────────────────────
\newcommand{\DefenseDate}{[JJ mois AAAA]}
\newcommand{\AcademicYear}{2025--2026}

% One \JuryMember{Name}{Position, institution}{Role} per line. Set \Jury to {}
% to remove the jury from the title page.
\newcommand{\Jury}{%
  \JuryMember{M. Prénom NOM}{Professeur, Université [X]}{Président du jury}
  \JuryMember{Mme Prénom NOM}{Maîtresse de conférences, Université [Y]}{Examinatrice}
  \JuryMember{M. Prénom NOM}{[Poste], [Laboratoire / Entreprise]}{Encadrant}
}

% ── Keywords ──────────────────────────────────────────────────────────────────
% Printed under the résumé/abstract and on the back cover.
\newcommand{\KeywordsFR}{mot-clé un, mot-clé deux, mot-clé trois}
\newcommand{\KeywordsEN}{keyword one, keyword two, keyword three}
